\uuid{pGFA}
\titre{Solution générale/condition initiale}
\auteur{hanani ; mzari}
\organisation{exo7}
\licence{CC BY-NC-SA 4.0}
\datecreate{2020-01-20}
\niveau{L1}
\module{Analyse}
\chapitre{Equation différentielle}
\sousChapitre{Résolution d'équation différentielle du premier ordre}
\difficulte{3}
\theme{Equations du premier ordre, solution générale, condition initiale}
\exo7qcmid{lille-2.equations-differentielles.010}

\contenu{
\begin{question}
\qclassification{225.01}
\qtags{motcle=solution générale/condition initiale}
On considère l'équation différentielle $(E)$ : $\displaystyle \sqrt{1+x^2}y'-xy=x$. Parmi les affirmations suivantes cocher celles qui sont vraies :
\begin{answers}
\bad{La solution générale de l'équation homogène est : $\displaystyle y=k\mathrm{e}^{1+x^2}$, $k\in \Rr$.}
\good{La solution générale de $(E)$ est : $\displaystyle y=-1+k\mathrm{e}^{\sqrt{1+x^2}}$, $k\in \Rr$.}
\good{Toute solution $y$ de $(E)$ vérifie $y'(0)=0$.}
\bad{$(E)$ admet une unique solution vérifiant $y'(0)=0$.}
\end{answers}
\begin{explanations}
La solution générale de l'équation homogène est $\displaystyle Y=k\mathrm{e}^{\sqrt{1+x^2}}$, $k\in \Rr$, et $y_0=-1$ est une solution particulière. En posant $x=0$ dans $(E)$, on obtient $y'(0)=0$, donc toute solution de $(E)$ vérifie $y'(0)=0$.
\end{explanations}
\end{question}
}
